% ============================================================
% Chapter 19 — Ethics & Privacy
% Source: PSYCON Engineering Design Document, Vol. 4, Ch. 19
% ============================================================
\section{Ethics and Privacy}
\label{sec:ethics-privacy}

\subsection{Ethical Principles}
\label{sec:ethical-principles}

\begin{itemize}
    \item Voluntary participation.
    \item Informed consent.
    \item Right to withdraw.
    \item Data minimization.
    \item Secure storage.
    \item Limited access.
    \item Transparent reporting.
\end{itemize}

\begin{requirementbox}
If recording speech, clearly explain what is collected, how it is processed, how long it is retained, and who can access it.
\end{requirementbox}

\subsection{Wearer Voice Profile}
\label{sec:wearer-voice-profile-privacy}

Speaker verification is biometric processing and requires consent separate from permission to process the conversation recording. Week~3 supports one local wearer profile created from three quality-checked recordings. Only the normalized averaged voice embedding and non-identifying enrollment metadata may be retained; the source enrollment recordings must be discarded immediately after processing. The profile must be encrypted using an operator-provisioned secret, excluded from logs and API responses, and replaceable or deletable by the operator.

The profile may identify only the enrolled wearer. Other voices remain anonymous recording-local speaker labels, and an uncertain comparison must abstain rather than select the nearest voice. Production deployment additionally requires authenticated access, encrypted transport, audit logging, a documented retention period, and measured false-match and false-rejection rates across the intended languages, microphones, and environments.
